%%%%%%%%%%%%%%%%%%%%%%%%%%%%%%%%%%%%%%%%
%% HiMCM/MidMCM LaTeX Template -- 2026 %%
%% Problem B: Requirement 1 working draft
%% Revised after reviewing user-supplied 2025 B papers, November 7, 2026.
%%%%%%%%%%%%%%%%%%%%%%%%%%%%%%%%%%%%%%%%
\documentclass[12pt]{article}
\usepackage{geometry}
\geometry{left=1in,right=0.75in,top=1in,bottom=1in}
\newcommand{\Problem}{B}
\newcommand{\Team}{1111111} % Replace with the actual Team Control Number.
\usepackage{newtxtext}
\usepackage{amsmath,amssymb,amsthm}
\usepackage{newtxmath}
\usepackage[pdftex]{graphicx}
\usepackage{xcolor}
\usepackage{fancyhdr}
\usepackage{microtype}
\usepackage{xurl}
\usepackage{needspace}
\usepackage{booktabs,tabularx,array}
\usepackage[hidelinks]{hyperref}
\setlength{\headheight}{15pt}
\setlength{\parskip}{3pt}
\setlength{\parindent}{1.2em}
\lhead{Team \Team}
\rhead{}
\cfoot{}
\newtheorem{theorem}{Theorem}
\newtheorem{corollary}[theorem]{Corollary}
\newtheorem{lemma}[theorem]{Lemma}
\newtheorem{definition}{Definition}

\begin{document}
\graphicspath{{.}}
\DeclareGraphicsExtensions{.pdf,.jpg,.png}
\thispagestyle{empty}
\vspace*{-16ex}
\centerline{\begin{tabular}{*3{c}}
\parbox[t]{0.3\linewidth}{\begin{center}\textbf{Problem Chosen}\\ \Large \Problem\end{center}}&
\parbox[t]{0.3\linewidth}{\begin{center}\textbf{2026\\ HiMCM/MidMCM\\ Summary Sheet}\end{center}}&
\parbox[t]{0.3\linewidth}{\begin{center}\textbf{Team Control Number}\\ \Large \textcolor{red}{\Team}\end{center}}\\
\hline
\end{tabular}}

\begin{center}
\vspace{1.5em}
{\Large\bfseries Measuring Useful Cooling in Open Systems}\\[0.8em]
{\bfseries Working summary: Requirement 1 only}
\end{center}
An open freezer case must maintain appropriate product temperatures despite heat exchange with its surroundings. Electricity use alone does not show how well it performs this task.

We distinguish thermal demand, electrical input, and delivered storage service. Infiltration, cabinet heat transfer, internal equipment, restocking, and defrost contribute to demand. Equipment condition modifies performance, while temperature differences across product locations determine whether the intended service is delivered.

We identify four quantities that define cooling service efficiency: energy input, service quantity, local service quality, and service duration. Compliant product mass-hours per kilowatt-hour provide one possible freezer measure. Product-specific temperature requirements remain a separate adequacy check. Comparisons require prescribed loading, matched ambient conditions, and complete operating cycles.

A correspondence table transfers these quantities to public spaces, where useful service becomes thermal-discomfort reduction accumulated over people and waiting time. Local exposure and a common baseline guide later comparisons of fans and evaporative coolers.

\medskip
\noindent\textit{Draft status.} Only Requirement 1 is developed here. Requirements 2--5 and the final whole-paper summary remain to be completed. No empirical performance results are claimed.

\clearpage
\pagestyle{fancy}
\rhead{Page \thepage}
\tableofcontents
\clearpage

\section{Understanding Cooling Efficiency}
\label{sec:understand}

An open freezer case is successful when it preserves the products placed inside it. Electricity use alone cannot establish this success: a case may consume less because it stores less food, operates for fewer hours, or fails to maintain the required temperatures. Conversely, high heat removal may reflect unwanted heat entering the opening. Cooling efficiency must therefore relate energy input to the service actually delivered.

\subsection{Selecting the quantities to evaluate}

We retain four quantities that have a direct role in both refrigeration and public-space cooling. Table~\ref{tab:service} states their meanings; each quantity enters either the useful-service numerator or the energy denominator.

\begin{table}[!ht]
\centering
\caption{Evaluation quantities and their transfer to a bus stop.}
\label{tab:service}
\small
\renewcommand{\arraystretch}{1.18}
\begin{tabularx}{\linewidth}{@{}>{\raggedright\arraybackslash}p{0.19\linewidth}
>{\raggedright\arraybackslash}X >{\raggedright\arraybackslash}X@{}}
\toprule
Quantity & Open freezer case & Bus-stop shelter\\
\midrule
Energy input & Electricity used by the case and its attributable refrigeration plant (kWh). & Electricity used by fans, pumps, and other cooling equipment (kWh).\\
Service quantity & Actual mass of products stored under a specified loading protocol (kg). & Number of people present in each occupied zone.\\
Local service quality & Product temperatures relative to product-specific storage requirements. & Thermal discomfort relative to the same setting without active cooling.\\
Service duration & Duration of product storage during the evaluation period (h). & Duration of passenger exposure while waiting (h).\\
\bottomrule
\end{tabularx}
\end{table}

\noindent\textbf{Energy input} establishes the resource cost. It must include the relevant auxiliaries and a consistent allocation of any shared plant, without overlapping meter readings. \textbf{Service quantity} distinguishes a small workload from a large one; declared loading limits prevent an overloaded cabinet from receiving artificial credit.

\textbf{Local service quality} establishes whether the task is achieved. Product temperatures near the opening and on different shelves can differ; measurements should represent the food rather than only the discharge air \cite{ashrae}. An average can hide a warm location, while uniform temperature can still be uniformly unsuitable. Product-temperature requirements must therefore be retained alongside energy performance \cite{etl}.

\textbf{Service duration} distinguishes a sustained outcome from a favorable snapshot. The evaluation period should include representative restocking, compressor cycling, and defrost recovery. Comparable initial and final thermal states prevent a short trial from gaining credit by merely using cold stored before the trial.

\clearpage
\subsection{Factors that explain performance}

The environment, cabinet, and operation explain changes in service and energy use. They enter the model as influencing factors rather than additional scores.

\noindent\textbf{Environmental exchange.} Air infiltration carries sensible heat and moisture through the opening. Cabinet walls transmit heat, and exposed surfaces exchange radiation with their surroundings. Opening geometry, insulation, ambient temperature, humidity, and air movement therefore affect the demand placed on the cooling system \cite{ashrae}.

\noindent\textbf{Equipment and operation.} Internal lights and fans release heat; warmer incoming stock adds a temporary load; defrost adds heat and interrupts normal cooling \cite{ashrae}. These processes explain variations in energy use and product temperature over a full cycle. Equipment age is relevant through actual condition---such as insulation deterioration, restricted airflow, or changed power demand---rather than through a separate age penalty.

\noindent\textbf{Spatial delivery.} Product arrangement and airflow determine where cooling is received. Evaluating each occupied location against its requirements captures this variation without an additional temperature-uniformity score.

\subsection{Implications for the bus-stop model}

The preceding analysis motivates a common evaluation principle:
\begin{equation}
\boxed{\text{Cooling efficiency}
=\frac{\text{useful service accumulated over objects and time}}
{\text{energy input over the same period}}.}
\label{eq:principle}
\end{equation}

For a freezer, accumulating compliant product mass over time gives a possible measure in kg\,h per kWh. At a bus stop, we instead accumulate thermal-discomfort reduction over people and waiting time. The next model will specify this response and spatial calculation. The two settings share a structure, not numerically comparable units.

Three adaptations are necessary. First, \textbf{evaluate people directly}: fans can improve the thermal sensation of an occupant without lowering room air temperature \cite{energystar}, so air-temperature reduction alone cannot fairly compare a fan with an evaporative cooler. The next model must account for local air speed, humidity, and radiation as well as temperature.

Second, \textbf{weight effects by occupancy}: cooling an empty area receives no passenger-service credit. Shelter geometry and cooled area determine which occupied zones benefit; area should not be multiplied by total headcount again. Outdoor air exchange can disperse cooled air while also changing the air movement experienced by passengers, so its effect must be evaluated through local conditions.

Third, \textbf{retain a minimum-service check}. In the freezer, individual product-temperature excursions must be checked even when aggregate service is high. At the bus stop, the model should recognize partial relief while also reporting remaining unacceptable heat exposure. Service adequacy is a separate requirement, rather than an extra weighted term that duplicates the efficiency numerator.

Within a venue, compare strategies under matched weather, occupancy, geometry, and evaluation periods. Water use, operating cost, and broader visitor experience enter the later effectiveness model. AI-assisted analysis and drafting are recorded in the appended report \cite{codex}.

\clearpage
\phantomsection
\addcontentsline{toc}{section}{References}
\begin{thebibliography}{9}
\bibitem{ashrae}
ASHRAE. \emph{2026 ASHRAE Handbook---Refrigeration},
Chapter 15, \emph{Retail Food Store Refrigeration and Equipment}.
Product Temperatures and Display Refrigerator Cooling Load sections.
\url{https://handbook.ashrae.org/Handbooks/R26/IP/r26_ch15/r26_ch15_ip.aspx}.
Accessed November 7, 2026.

\bibitem{etl}
UK Department for Energy Security and Net Zero.
\emph{Energy Technology List: Retrofit Measures for Open Fronted Refrigerated Display Cabinets}.
Definitions and product-temperature performance requirements.
\url{https://etl.energysecurity.gov.uk/products/refrigeration-equipment/retrofit-measures-open-fronted-refrigerated-display-cabinets}.
Accessed November 7, 2026.

\bibitem{energystar}
ENERGY STAR. \emph{Ceiling Fan Installation and Usage Tips}.
\url{https://www.energystar.gov/products/ceiling_fans/installation-and-usage-tips}.
Accessed November 7, 2026.

\bibitem{codex}
OpenAI. \emph{Codex}, AI-assisted formulation, reference checking,
English drafting, and LaTeX preparation, November 2026.
Exact deployed model identifier to be recorded from the session interface.
See the preliminary AI-use record following this draft.

\bibitem{aipolicy}
COMAP. \emph{Use of AI Tools in COMAP Contests}, v102025.
\url{https://www.contest.comap.com/undergraduate/contests/mcm/flyer/Contest_AI_Policy.pdf}.
Accessed November 7, 2026.
\end{thebibliography}

% Add Requirements 2--5, the memorandum, and the blog before References.
% Replace the working summary after completing the full analysis.
\clearpage
\phantomsection
\addcontentsline{toc}{section}{Report on Use of AI Tools}
\section*{Report on Use of AI Tools}
\noindent\textbf{Preliminary record for this partial draft.}

OpenAI Codex assisted with problem translation, the overall modeling discussion,
and the development of Requirement 1 in November 2026. The exact deployed model
identifier must be recorded from the session interface.

For this edit, the user requested an English section in the supplied template
covering infiltration, cabinet heat transfer, internal equipment, restocking,
defrost and operating cycles, equipment age, and spatial temperatures. This is
a summary of the request, not a verbatim transcript.

The assistant proposed and subsequently simplified the evaluation framework,
checked source material, drafted the English text, and prepared the LaTeX.
Separate AI review passes checked technical consistency. At the user's request,
the 2025 Problem B solutions of Teams 16442, 16738, and 16739 were reviewed as
structural references. Their numerical estimates, weighting procedures, and
wording were not used in the present model. This revision emphasizes indicator
justification and transfer to the later task. No field measurements or empirical
performance results were produced.

Before submission, add the exact relevant prompts and complete outputs,
including earlier retained contributions and any additional AI tools, as
required by the contest reporting guidance \cite{aipolicy}. Retain this source
file and the chat as records. The student team's independent verification is
still required and has not been asserted in this preliminary record.

\end{document}

